% Generated draft: requires mathematical and source review.
\documentclass[11pt]{article}
\usepackage[margin=25mm]{geometry}
\usepackage{fontspec}
\setmainfont{Latin Modern Roman}
\usepackage{amsmath,amssymb,mathtools,unicode-math}
\setmathfont{Latin Modern Math}
\usepackage{enumitem,needspace,xurl,hyperref}
\setlength{\parindent}{0pt}
\setlength{\parskip}{5pt}
\setlength{\emergencystretch}{3em}
\newcommand{\N}{\mathbb{N}}
\newcommand{\Z}{\mathbb{Z}}
\newcommand{\Q}{\mathbb{Q}}
\newcommand{\R}{\mathbb{R}}
\newcommand{\C}{\mathbb{C}}
\newcommand{\F}{\mathbb{F}}
\newcommand{\A}{\mathbb{A}}
\newcommand{\PP}{\mathbb{P}}
\newcommand{\E}{\mathbb{E}}
\newcommand{\eps}{\varepsilon}
\newcommand{\dd}{\,\mathrm{d}}
\begin{document}
\section*{UnsolvedMath ID: 6700075}
\textbf{Source problem code:} catalogue entry
\subsection{Definitions and mathematical statement}
For $Y\subset W$, the source uses the remnant-of-collapse scalar-curvature quantity $\operatorname{Sc}^{[n]}_{g,\cap}(Y)$ but does not supply its definition. Retain the predicate
\[
P_{g,n}(Y):\quad\operatorname{Sc}^{[n]}_{g,\cap}(Y)=+\infty.
\]
It conjectures three properties:
\begin{enumerate}
\item $P_{g,n}(Y)$ is independent of the ambient Riemannian metric $g$ on the fixed $W$.
\item The predicate is invariant under bi-Lipschitz equivalence.
\item The predicate is not invariant under homeomorphism.
\end{enumerate}
The bi-Lipschitz assertion requires a specified domain: a bi-Lipschitz bijection controls distances in both directions, but the source does not say whether it is between the sets alone or between ambient pairs $(W,Y)$. Non-topological-invariance similarly requires examples of the intended homeomorphic objects with different predicate values.

The source explicitly says there is no serious evidence and offers heuristic examples. For $1\leq k\leq n-2$ and $\mu>0$, it describes nested neighborhoods $V_{\varepsilon_i}$ of flat $k$-subtori in $W=\mathbb T^{n+1}$, with boundary scalar curvatures tending to $+\infty$ and intersection $Y$ of Lebesgue measure $\operatorname{mes}(W)-\mu$. Necessarily $\mu\leq\operatorname{mes}(W)$ for this measure value to be nonnegative; that restriction is absent from the stated range. The meaning of curvature tending to infinity on varying boundaries is not specified. It also suggests making all neighborhoods homeomorphic to a $k$-ball, which conflicts dimensionally with full-dimensional neighborhoods when $k<n+1$ unless another object is intended. Finally it tentatively proposes compact zero-dimensional Cantor sets $Y\subset\mathbb R^{n+1}$ with $\operatorname{Sc}^{[n]}_{\cap}(Y)\ne+\infty$. None of these examples is promoted here to a verified construction.
\subsection{Short English statement}
The source conjectures that infinite remnant-of-collapse scalar curvature is independent of the ambient smooth metric and invariant under bi-Lipschitz equivalence, but not under homeomorphism. Its invariant definition and several details of the heuristic examples are missing or inconsistent.
\subsection{Sources}
\begin{itemize}
\item Misha Gromov, \emph{101 Questions, Problems and Conjectures around Scalar Curvature} (2017), printed tag [?79], occurrence 78, PDF page 78, \url{https://www.ihes.fr/~gromov/wp-content/uploads/2018/08/101-problemsOct1-2017.pdf}; accessed 2026-07-29, PDF-text-normalized extraction. The provenance describes the public author PDF as incomplete.
\item The Cantor-set comparison cites section 5.3 of source reference [48], whose bibliography entry is not supplied.
\item Local registry \texttt{data/unsolvedmath/problems.json}, permanent UnsolvedMath ID $6700075$, source code catalogue entry, \emph{Scalar Curvature Question [?79]}. Dataset attribution: Ulam AI and the UnsolvedMath Contributors.
\end{itemize}
\end{document}
